\[f(x) = a \ln (x).\]

On note $\mathcal{C}$ sa courbe représentative dans un repère orthonormé.

Soit $x_{0}$ un réel strictement supérieur à 1.

\medskip

\begin{enumerate}
\item 
On a $f(x) = 0 \iff a\ln x = 0 \iff \ln x = 0$ (car $a \ne 0$) et par croissance de la fonction exponentielle $\e^{\ln x} = \e^0 \iff x = 1$.
\item 
$F$ est une différence de fonctions dérivables sur $]0~;~+\infty[$, donc sur cet intervalle :

$F'(x) = a\left[\ln (x) + x \times \dfrac 1x - 1\right] = a[\ln (x) + 1 - 1] = a\ln (x) = f(x)$, ce qui montre que $F$ est une primitive de $f$ sur $]0~;~+\infty[$.
\item 
On a vu que $\mathcal{C}$ coupe l'axe des abscisses en $x = 1$ donc la surface bleue correspond à des points où $x \geqslant 1$, soit $\ln (x) \geqslant 0 \Rightarrow a\ln (x) \geqslant 0$.

Autrement dit  pour $x \geqslant 1$, la fonction $f$ est positive et on sait que sur un intervalle $[1~;~x_0]$ avec $x_0 \geqslant 1$, l'aire de la surface limitée par sa courbe représentative, l'axe des abscisses et les droites d'équations $x = 1$ et $x = x_0$ est égale à l'intégrale :
\[\displaystyle\int_1^{x_0} f(x)\:\text{d}x = [F(x)]_1^{x_0} = F(x_0) - F(1) = a[x_0 \ln (x_0) - x_0] - a [1 \ln (1) - 1] =.\]

l'aire bleutée est en unités d'aire : $ a [x_0 \ln (x_0) - x_0] + a = a\left[x_0 \ln (x_0) - x_0 + 1\right]$.
\end{enumerate}

\begin{center}
\psset{unit=2cm,arrowsize=2pt 3}
\begin{pspicture*}(-1,-1)(4.5,2)
\psaxes[linewidth=1.25pt,Dx=5,Dy=5]{->}(0,0)(-1,-1)(4.5,2)
\pscustom[fillstyle=solid,fillcolor=blue]{\psplot[plotpoints=2000,linewidth=1.25pt,linecolor=blue]{1}{1.8}{x ln 1.3 mul}\psline(1.8,0.764)(1.8,0)(1,0)}
\psplot[plotpoints=2000,linewidth=1.25pt,linecolor=blue]{0.1}{4.5}{x ln 1.3 mul}
\uput[dr](1,0){1}\uput[d](1.8,0){$x_0$}\uput[ul](1.8,0.764){$M$}
\uput[dr](4,1.6){\blue $\mathcal{C}$}\uput[dl](0,0){O}
\uput[d](4.4,0){$x$}\uput[l](0,1.9){$y$}
\end{pspicture*}
\end{center}

On note $T$ la tangente à la courbe $\mathcal{C}$ au point $M$ d'abscisse $x_{0}$.

On appelle $A$ le point d'intersection de la tangente $T$ avec l'axe des ordonnées et $B$ le projeté orthogonal de $M$ sur l'axe des ordonnées.

\begin{center}
\psset{unit=2cm,arrowsize=2pt 3}
\begin{pspicture*}(-1,-1)(4.5,2)
\psaxes[linewidth=1.25pt,Dx=5,Dy=5]{->}(0,0)(-1,-1)(4.5,2)
\psplot[plotpoints=2000,linewidth=1.25pt,linecolor=blue]{0.1}{4.5}{x ln 1.3 mul}
\uput[dr](1,0){1}\uput[d](1.8,0){$x_0$}\uput[ul](1.8,0.764){$M$}
\uput[dr](4,1.6){\blue $\mathcal{C}$}\uput[dl](0,0){O}
\psline[linestyle=dashed](1.8,0.764)(1.8,0)
\uput[d](4.4,0){$x$}\uput[l](0,1.9){$y$}
\psplotTangent{1.8}{4}{x ln 1.3 mul}
\psline(-1,0.764)(4.5,0.764)\uput[ul](0,0.764){B}\uput[ul](0, -0.5){A}
\psframe(0,0.764)(0.1,0.664)\uput[ul](3.1,1.8){$T$}
\psline[linewidth=1.5pt,linecolor=red](0,-0.5)(0,0.764)\uput[l](0,0.15){\red $a$}
\psline[linewidth=1.5pt,linecolor=red](2.71828,0)(2.71828,1.3)\uput[r](2.71828,0.65){\red $a$}
\uput[d](2.71828,0){e}
\end{pspicture*}
\end{center}

\begin{enumerate}
\setcounter{enumi}{3}
\item 


On sait (équation de la tangente au point d'abscisse $x_0$) que :

$M(x~;~y)\in T \iff y - f(x_0) = f'\left(x - x_0) \right)$.

$\bullet~~$$f\left(x_0\right) = a\ln x_0$ ;

$\bullet~~$$f$ est dérivable sur $]0~;~+ \infty[$ et sur cet intervalle $f'(x) = \dfrac ax$, donc $f'\left(x_0 \right) = \dfrac{a}{x_0}$.

On obtient donc :

$M(x~;~y)\in T \iff y - a\ln x_0 = \dfrac{a}{x_0} \times \left(x - x_0 \right) \iff y = a\ln x_0 + \dfrac{ax}{x_0} - a$.

En particulier $T$ coupe l'axe des ordonnées si $x = 0$, d'où $y = a\ln x_0 - a$ (ordonnée de A).

L'ordonnée de B est égale à $f\left(x_0\right) = a\ln \left(x_0\right)$.

On a $\text{AB} = \left|y_{\text{B}} - y_{\text{A}} \right| = \left|a\ln \left(x_0\right) - \left(a\ln x_0 - a\right) \right| = \left|a \right| = a$ (car $a > 0$).

\emph{Remarque} : On a $f(\e) = a \ln \e = a \times 1 = a$.

$f(\e) = a$ : on a mis en évidence ceci sur le dessin ; le corollaire est que la tangente à la courbe au point d'abscisse e contient l'origine O !
\end{enumerate}
